\section{The earlier four-term reduction}\label{sec:prior}
For a primitive reduced quadratic form $[a,b,c]$ of discriminant $-D$, the corresponding conjugate has modulus within $2079$ of $\exp(\pi\sqrt D/a)$ \cite{riffaut}. This gives an archimedean complement to the disk lemma. If $t$ conjugates of increasing moduli $r_i$ have successive ratios at least 16, every generalized Vandermonde minor on them and $t$ distinct exponents $n_1<\cdots<n_t$ is nonzero. Scale column $i+1$ relative to column $i$ by $(r_ir_{i+1})^{-(n_{i+1}-n_i)/2}$. Moving away from the diagonal reduces entries by at least a factor four at every step. Each row's off-diagonal sum is at most $2\sum_{k\ge1}4^{-k}=2/3$ times its diagonal magnitude. Strict diagonal dominance proves the assertion.

For consecutive leading coefficients $a<b$ the integer inequalities
\[
       D\ge16a^2,\qquad 9D(b-a)^2\ge16a^2b^2                         \tag{2}
\]
suffice for a ratio greater than 16. Indeed the first exponential is at least $e^{4\pi}>34\cdot2079$ and the exponential ratio is at least $e^{4\pi/3}>32$, so the two cusp errors preserve the claimed separation.

\heading{What was inherited and verified.} The incoming CM report \cite{prior} applied the quadratic factorization $\zeta_K=\zeta L(\cdot,\chi_{d_K})$ and Greni\'e--Molteni's Corollary 1.2 \cite{gm} for $D\ge10^{11}$. It obtained an odd split prime coprime to the conductor in each of $(0,.046R]$, $(.048R,.189R]$, $(.230R,.500R]$ for $R=\sqrt D$. With the principal form these give four separated conjugates. We reproduced its exact endpoint checks, yielding margins $68,191,751$. Its finite sieves gave the following counts; we reran both.
\begin{center}\small
\begin{tabular}{@{}p{.26\linewidth}p{.65\linewidth}@{}}\toprule
Range or step & Outcome \\ \midrule
$10^7\le D<10^{11}$ &
Of $49{,}995{,}000{,}000$ order discriminants, $12{,}498{,}750{,}000$ pass at 2, $37{,}449{,}133{,}804$ via powers of an odd split prime and $47{,}116{,}196$ via three individually checked primes. \\[2pt]
$3\le D\le10^7$ &
Of $5{,}000{,}000$ discriminants, $67$ have $h<4$. Integer separation certifies $4{,}999{,}603$. Of the remaining $330$, interval separation removes $228$ and a determinant comparison removes $22$. \\[2pt]
Residual $80$ & $50$ have degree four, $29$ degree five, and $-4075$ has degree six. The earlier report covers $-4075$ by complex intervals. \\ \bottomrule
\end{tabular}
\end{center}
The disk proof gives more at $-4075$. At 7 all six conjugates are in one nonzero residue disk of depth at least one, by Section~\ref{sec:growth}. The proof of (1) for $p\ge t$ also works for $t=6$, giving threshold $1/2$ at 7. Hence $s(-4075)=7$. The following sections close the 79 earlier cases and provide a separate unconditional global proof through the stronger five-term reduction.
